\documentclass[10pt]{article}
\usepackage[T1]{fontenc}
\usepackage[utf8]{inputenc}
\usepackage[french]{babel}
\usepackage{amsmath,amssymb,mathrsfs,mathpazo}
\usepackage{geometry}
\geometry{a4paper,margin=2cm}
\usepackage{graphicx,array,tabularx,tabularray,colortbl,multirow,multicol,booktabs,enumitem,pifont,ulem,fancybox,diagbox,cancel}
\usepackage{pstricks,pstricks-add,pst-plot,pst-tree,pst-node,pst-text,pst-func,pst-math,pst-bspline,pst-3dplot,pst-eucl,pst-solides3d}
\let\pscalculate\undefined
\usepackage{pst-bezier}
\usepackage{tikz,pgfplots,tkz-tab}
\usetikzlibrary{arrows,arrows.meta,patterns,calc,positioning,trees,decorations.markings,shapes.geometric,shapes.misc,shapes.arrows,intersections,decorations.pathmorphing,decorations.pathreplacing,backgrounds,shadows,babel}
\pgfplotsset{compat=1.18}
\usepgfplotslibrary{fillbetween}
\usepackage[np]{numprint}
\usepackage{listings,makecell,pas-tableur,twoopt}
\let\barre\undefined
\usepackage[pstricks]{tablvar}
\definecolor{darkgreen}{RGB}{0,100,0}
\providecommand{\red}{\color{red}}
\providecommand{\blue}{\color{blue}}
\providecommand{\green}{\color{green}}
\providecommand{\black}{\color{black}}
\providecommand{\white}{\color{white}}
\providecommand{\gray}{\color{gray}}
\providecommand{\cyan}{\color{cyan}}
\providecommand{\magenta}{\color{magenta}}
\providecommand{\yellow}{\color{yellow}}
\providecommand{\violet}{\color{violet}}
\providecommand{\orange}{\color{orange}}
\providecommand{\euro}{\texteuro}
\providecommand{\EUR}{\texteuro}
\providecommand{\ieme}{\textsuperscript{e}}
\providecommand{\degres}{\ensuremath{^\circ}}
\providecommand{\degree}{\ensuremath{^\circ}}
\providecommand{\vv}[1]{\overrightarrow{#1}}
\providecommand{\Coord}[1]{#1}
\providecommand{\decosix}{}
\providecommand{\starredbullet}{\ensuremath{\bullet}}
\pagestyle{empty}
\setlength{\parindent}{0pt}
\usepackage{adjustbox}\DeclareUnicodeCharacter{2236}{:}\DeclareUnicodeCharacter{2190}{\ensuremath{\leftarrow}}
\providecommand{\R}{\mathbb{R}}
\providecommand{\N}{\mathbb{N}}
\providecommand{\Z}{\mathbb{Z}}
\providecommand{\Q}{\mathbb{Q}}
\providecommand{\D}{\mathbb{D}}
\providecommand{\C}{\mathbb{C}}
\providecommand{\e}{{\mathrm{e}}}
\providecommand{\vect}[1]{\overrightarrow{#1}}
\providecommand{\vectt}[1]{\overrightarrow{\mathrm{#1}}}
\providecommand{\vvt}[1]{\overrightarrow{\mathrm{#1}}}
\providecommand{\np}[2][]{{#2~#1}}
\providecommand{\vv}[1]{{\overrightarrow{#1}}}
\renewcommand{\d}{{\mathrm{d}}}
\providecommand{\degres}{{^\circ}}
\providecommand{\degre}{{^\circ}}
\providecommand{\blue}{\color{blue}}
\providecommand{\red}{\color{red}}
\providecommand{\enskip}{\quad}
\providecommand{\llbracket}{\lbrack\!\lbrack}
\providecommand{\rrbracket}{\rbrack\!\rbrack}
\providecommand{\euro}{{\text{€}}}
\providecommand{\no}{{\text{n°}}}
\providecommand{\interval}[2]{{[#1;#2]}}
\providecommand{\intervFO}[2]{{[#1;#2[}}
\providecommand{\dotfill}{\ldots}

\providecommand{\ds}{\displaystyle}
\providecommand{\textding}[1]{#1}
\providecommand{\Oij}{\(\mathrm{O};\vec{\imath},\vec{\jmath}\)}
\providecommand{\Oijk}{\(\mathrm{O};\vec{\imath},\vec{\jmath},\vec{k}\)}
\providecommand{\Ouv}{\(\mathrm{O};\vec{u},\vec{v}\)}
\providecommand{\intervFF}[2]{\ensuremath{[#1;#2]}}
\renewcommand{\euro}{\ifmmode\text{\texteuro}\else\texteuro\fi}
\DeclareUnicodeCharacter{00B0}{\ifmmode^{\circ}\else\textdegree\fi}
\DeclareUnicodeCharacter{20AC}{\ifmmode\text{\texteuro}\else\texteuro\fi}
% Énoncés : Éducation nationale. Transcription et corrigés : APMEP, Denis Vergès et contributeurs. Découpage et classement : CoopMaths, Sébastien Lozano, Éric Elter et contributeurs. Adaptation de la présentation : Maths Explorer.
% CC BY-SA : https://coopmaths.fr/www/about/
% https://coopmaths.fr/alea/static/eam/2026/tex/eam_specialite_2026_06_ameriquenord_2.tex
\begin{document}
\noindent\textit{APMEP / CoopMaths. CC BY-SA. Présentation adaptée par Maths Explorer.}\par\medskip
Pour réduire sa facture d'électricité, Camille a décidé de faire poser des panneaux solaires sur le toit de sa maison.

Elle souhaite analyser sa production et estimer le temps nécessaire pour rentabiliser cet investissement.

\textbf{Les deux parties suivantes sont indépendantes.}

\subsection*{Partie A}

Lors d'une belle journée ensoleillée, la puissance électrique en kilowatt (kW) des panneaux solaires de Camille peut être modélisée en fonction de l'heure par une fonction $f$. On admet que $f$ est définie sur $[0 ; 24]$ et on donne sa courbe représentative $\mathcal{C}_{f}$ ci-dessous.

Avec la précision permise par le graphique :

\begin{center}
	\begin{tikzpicture}
		\begin{axis}[/pgf/number format/.cd, use comma,
			x={6mm}, y={6mm},
			xmin = -1, xmax = 24.5,
			ymin = -0.9, ymax = 10.6,
			xtick distance=1,
			ytick distance=1,
			minor tick num = 1,
			grid=both, axis lines =center,
			grid style={gray!,line width=0.35pt,dotted}]
			\addplot[line width=1.2pt,color=red,smooth,samples=100,domain= 0:24] {8*exp(-0.08*(x-13)*(x-13))};
			\node[below left=1mm] at (axis cs: 0,0) {{\footnotesize 0}};
			\node at(axis cs: 8,2) {$ \mathcal{C}_f $};
			\node[text width=2cm] at(axis cs: 2,9.25) {Puissance électrique (en kW)};
			\node at(axis cs: 21.6,1) {Horaire de la journée};
		\end{axis}
	\end{tikzpicture}
\end{center}















\begin{enumerate}
	\item Donner la puissance électrique des panneaux solaires à 11 h 00.

	\item Résoudre graphiquement l'inéquation \quad $f(x) \geqslant 5$ \quad et interpréter ce résultat dans le contexte de l'énoncé.
\end{enumerate}

\subsection*{Partie B}
Le coût pour 1 kilowattheure (kWh) consommé au tarif réglementé était de 0,15~\euro{} en 2020. On admet que ce tarif réglementé augmente de 6~\% chaque année.

On note $c_{n}$ le coût en euros (\euro{}) pour 1 kWh consommé durant l'année $2020+n$, avec $n$ un entier naturel. On a alors $c_{0}=0,15$.

\begin{enumerate}
	\item Déterminer la nature de la suite $\left(c_{n}\right)$. On précisera sa raison.

	\item Pour tout entier naturel $n$, exprimer $c_{n}$ en fonction de $n$.

	\item Donner le calcul permettant d'obtenir le coût pour 1 kWh consommé en 2030. II n'est pas demandé d'effectuer ce calcul.

	\item On admet que, chaque année depuis 2020, l'utilisation des panneaux solaires de Camille lui a permis d'éviter l'achat de \np[kWh]{2000} par an.
	L'installation des panneaux solaires en janvier 2020 a coûté à Camille 7000~\euro{}.
	On considère le programme Python ci-contre.

	\begin{minipage}{0.5\linewidth}
		\begin{enumerate}
			\item Dans le contexte de l'énoncé, que représentent les variables \texttt{c} et \texttt{S} du programme ?

			\item On exécute le programme ci-contre.

			II affiche 16.

			Interpréter ce résultat dans le contexte de l'énoncé.
		\end{enumerate}
	\end{minipage}
\hfill
\begin{minipage}{0.37\linewidth}
\hfill
\begin{ttfamily}
\begin{tblr}{hline{1,Z}={solid}, vlines,colspec={X},rowsep=0pt}
	n = 0\\
	c = 0.15\\
	S = 0\\
	while S < 7000 :\\
	\qquad S = S + c*2000\\
	\qquad n = n+1\\
	\qquad c = 1.06* c\\
	print(n)\\
\end{tblr}
\end{ttfamily}
\end{minipage}
\hfill~
\end{enumerate}

\end{document}
